\documentclass{article}
\usepackage[utf8]{inputenc}
\usepackage{geometry}
\geometry{a4paper, margin=1in}

\title{Optimism and the Power of Closing the Hypocrisy Gap}
\author{Albert Jan van Hoek}
\date{\today}

\begin{document}

\maketitle

I remember how excited I once was about Elon Musk and the early days of Tesla. The future felt bright.

I also remember enjoying saying ``You are fired!'' Trump \& The Apprentice was grotesque — but it was about efficiency and confidence. About doing something. Changing something. Being rewarded.

\section*{Why I’m Optimistic}

People ask why I’m still optimistic.

I am optimistic. Not because powerful figures will save us. Not because politics will suddenly become noble. Not because billionaires will behave differently.

I’m optimistic for another reason.

In my \textbf{Theory of Long-term Collaboration}, stable relationships are created — or broken — by the alignment between what we say and what we do. When words and behaviour reinforce each other, cooperation becomes self-strengthening. When they diverge, systems erode.

\section*{Hypocrisy as Potential}

Seen like this, hypocrisy is our real enemy.

Hypocrisy is the gap between knowing and living.

We know what the other needs.
We know what our body needs.
We know what society needs.
We know what the planet needs.

Yet we hesitate. We don’t provide it.

But hypocrisy is not the end. It is the beginning.
It is stored potential.
Tension before movement.

You can only be hypocritical about something you already understand. The insight is there. The behaviour is not — yet.

\section*{The Path Forward}

We live in a deeply interconnected and interdependent world. Our behaviour toward each other makes the difference. Fixing small personal hypocrisies already shifts the system. When we do it together, the effect becomes exponential.

We are wired to collaborate — as biological and learning systems. Honesty, trust and forgiveness are energetically cheaper than anxiety, stress and distrust. Alignment requires less energy than internal conflict.

Sometimes we cling to structures that feel stabilizing but prevent adjustment. Money is one of them. Money is a tool for coordination — powerful, useful — but survival at its core depends on food, shelter, safety, care, trust and cooperation. Not on money itself. When we equate money with survival, fear grows. And fear widens the gap between what we know and how we live.

\section*{A Call to Action}

That is why I remain optimistic.

The problem is not far away. It is within reach.

Overcoming our own hypocrisy is a manageable task.

Therefore LinkedIn — I’m off for a run.
One hypocrisy out of the window at a time.
Running toward a better place.

Just do it!

\end{document}
